\section{Objetivos y Preguntas de Investigación}

\subsection{Preguntas de Investigación}
\textbf{Pregunta General:}
¿Cómo se relacionan los ingresos de los hogares con el costo de una canasta básica de referencia en Chile, Brasil y Argentina, según los microdatos y períodos disponibles para cada país?

\textbf{Preguntas Específicas:}
\begin{enumerate}
    \item ¿Qué medidas de ingreso del hogar y de costo de una canasta básica de referencia pueden construirse con fuentes documentadas para cada país, período y cobertura geográfica?
    \item ¿Qué proporción de hogares puede cubrir el costo de su canasta básica local con el ingreso monetario disponible, y cómo varía descriptivamente según nivel de ingreso, sexo y tamaño del hogar?
    \item ¿Qué diferencias en las encuestas, los conceptos de ingreso, la cobertura territorial y las canastas oficiales limitan la comparación entre los tres países?
\end{enumerate}

\subsection{Objetivos}
\textbf{Objetivo General:}
Comparar descriptivamente la capacidad de los hogares de Chile, Brasil y Argentina para cubrir el costo de una canasta básica de referencia con sus ingresos monetarios, utilizando las fuentes y períodos disponibles y explicitando los límites de comparabilidad.

\textbf{Objetivos Específicos:}
\begin{enumerate}
    \item Identificar, documentar y armonizar con cautela los ingresos, ponderadores y variables sociodemográficas disponibles en CASEN, PNAD Contínua y EPH.
    \item Reunir para cada país el valor de la medida de costo de vida seleccionada, identificando si su fuente es gubernamental o independiente, así como unidad, mes, territorio y método.
    \item Estimar, cuando los insumos sean compatibles, la razón entre ingreso mensual y costo de la canasta en una unidad equivalente, además de la proporción ponderada de hogares bajo cobertura.
    \item Describir la heterogeneidad de los indicadores por grupos sociales y documentar las diferencias que impiden interpretar los resultados como una clasificación directa entre países.
\end{enumerate}